\chapter{Suffix Structures \& Document Similarity (Module 2)}
\label{ch:suffix_structures}

\section{Introduction to Suffix Structures}
While linear string matchers locate fixed patterns, full-text indexing requires answering arbitrary substring queries without re-scanning the underlying corpus. Suffix structures—namely Suffix Arrays, the Kasai LCP array, and Suffix Automata—preprocess a document of length $N$ to permit sub-linear query evaluation and cross-document similarity extraction.

\section{Suffix Arrays \& Prefix Doubling}
\subsection{Mathematical Definition}
Let $T$ be a text of length $N$ terminated by an implicit sentinel character $\$ < c \ \forall c \in \Sigma$. The Suffix Array $SA$ is a permutation of integers $\{0, \dots, N - 1\}$ representing the lexicographically sorted order of all suffixes of $T$:
\begin{equation}
T[SA[0] \dots N-1] < T[SA[1] \dots N-1] < \dots < T[SA[N-1] \dots N-1]
\end{equation}

\subsection{Prefix-Doubling Algorithm in $O(N \log^2 N)$}
Prefix-doubling sorts suffixes by their prefixes of length $2^k$ for $k = 0, 1, 2, \dots, \lceil \log_2 N \rceil$. In iteration $k$, the rank of suffix $i$ is determined as a 2-tuple:
\begin{equation}
\text{Tuple}(i) = \left( \text{rank}_{2^{k-1}}(i), \ \text{rank}_{2^{k-1}}(i + 2^{k-1}) \right)
\end{equation}
EduTrack sorts these tuples using a hand-built merge-sort (\texttt{customSort} in \texttt{SuffixArray.java}), strictly avoiding \texttt{java.util.Arrays.sort}. Once constructed, any query pattern $P$ of length $M$ is located via binary search on $SA$ in $O(M \log N)$ time.

\subsection{Short Illustrative Example (Independent from Project)}
Consider the canonical string $T = \text{"banana"}$ ($N = 6$):
\begin{table}[H]
\centering
\begin{tabular}{cclc}
\toprule
$i$ & $SA[i]$ & Suffix String & $LCP[i]$ \\
\midrule
0 & 5 & \texttt{"a"} & 0 \\
1 & 3 & \texttt{"ana"} & 1 \\
2 & 1 & \texttt{"anana"} & 3 \\
3 & 0 & \texttt{"banana"} & 0 \\
4 & 4 & \texttt{"na"} & 0 \\
5 & 2 & \texttt{"nana"} & 2 \\
\bottomrule
\end{tabular}
\caption{Suffix Array and LCP Values for \texttt{"banana"}}
\label{tab:banana_sa}
\end{table}

\subsection{EduTrack Domain Application: Textbook Corpus Indexing}
In EduTrack (\texttt{SuffixArray.java}), the Suffix Array indexes university textbooks and reference documents (e.g., \texttt{algorithms\_handbook.txt}). Binary search returns the exact text offset of technical terms in sub-millisecond time.

\subsection{Screenshot \& Verification Specifications}
\begin{figure}[H]
    \centering
    \fbox{\begin{minipage}{0.85\textwidth}
        \centering
        \vspace{1.2cm}
        \textbf{[PLACEHOLDER: Screenshot 5 - Suffix Array Binary Search]}\\[0.2cm]
        \textit{Capture the GUI/CLI executing Suffix Array binary search on library text.}\\
        \textit{Display: Search for term "Suffix", binary search iterations, and detected index offset.}
        \vspace{1.2cm}
    \end{minipage}}
    \caption{Suffix Array Construction and Binary Search in EduTrack}
    \label{fig:screenshot_sa}
\end{figure}

% ------------------------------------------------------------------------------
\section{Linear-Time Suffix Array Induced Sorting (SA-IS)}
\subsection{Theoretical Foundations \& S/L Suffix Classification}
The SA-IS algorithm~\cite{nong2009} constructs the suffix array in linear time $O(N)$ using induced sorting. Every suffix $i$ is classified as:
\begin{itemize}
    \item \textbf{S-type:} $T[i \dots N-1] < T[i+1 \dots N-1]$.
    \item \textbf{L-type:} $T[i \dots N-1] > T[i+1 \dots N-1]$.
\end{itemize}
A suffix $i$ is a \textbf{Left-Most S-type (LMS)} character if $T[i]$ is S-type while $T[i-1]$ is L-type.

\subsection{Two-Phase Induced Sorting Invariant}
\begin{enumerate}
    \item \textbf{Bucket Placement:} Buckets for each character are established with head and tail pointers. LMS characters are placed into the tails of their respective buckets.
    \item \textbf{Induce L-type Suffixes:} Scan $SA$ from left to right ($0$ to $N-1$). For each entry $SA[i]$, if $SA[i] - 1$ is L-type, insert it into the current head of bucket $T[SA[i] - 1]$.
    \item \textbf{Induce S-type Suffixes:} Scan $SA$ from right to left ($N-1$ down to $0$). For each entry $SA[i]$, if $SA[i] - 1$ is S-type, insert it into the current tail of bucket $T[SA[i] - 1]$.
\end{enumerate}
If LMS substrings are not pairwise unique, the problem recursively reduces on a compressed alphabet of size $\le N/2$. The recurrence $T(N) = T(N/2) + O(N)$ resolves to $O(N)$ worst-case time.

\subsection{EduTrack Implementation Details}
EduTrack implements SA-IS in \texttt{SAIS.java} entirely using primitive arrays (\texttt{int[]}, \texttt{boolean[]}) with zero object allocation overhead.

% ------------------------------------------------------------------------------
\section{Kasai's Longest Common Prefix (LCP) Algorithm}
\subsection{Mathematical Definition \& The Kasai Invariant}
The LCP array stores the length of the longest common prefix between consecutive suffixes in the sorted Suffix Array:
\begin{equation}
LCP[i] = \text{lcp}(T[SA[i] \dots N-1], \ T[SA[i-1] \dots N-1]), \quad LCP[0] = 0
\end{equation}
Let $rank[j]$ denote the position of suffix $j$ in $SA$. Kasai et al.~\cite{kasai2001} discovered a profound structural invariant:
\begin{theorem}[Kasai's Invariant]
Let $h$ denote the LCP between suffix $i$ and its predecessor in the suffix array, $k = SA[rank[i] - 1]$. When moving from suffix $i$ to suffix $i + 1$, the common prefix length with its new predecessor satisfies:
\begin{equation}
LCP[rank[i + 1]] \ge LCP[rank[i]] - 1
\end{equation}
\end{theorem}
\begin{proof}
Since suffix $i$ and suffix $k$ share a prefix of length $h$, suffix $i + 1$ and suffix $k + 1$ must share a prefix of length $h - 1$. In the sorted suffix array, suffix $k + 1$ precedes suffix $i + 1$. Because the LCP between any two suffixes is the minimum LCP along the intervening range in $SA$, the LCP between suffix $i + 1$ and its immediate predecessor cannot be less than $h - 1$.
\end{proof}
Because $h$ increments at most $N$ times and decrements at most $N$ times, the total character comparisons across the entire algorithm is bounded by $2N$, guaranteeing \textbf{deterministic $O(N)$ running time}.

\subsection{EduTrack Domain Application: Cross-Document Plagiarism Detection}
EduTrack leverages Kasai's algorithm (\texttt{KasaiLCP.java}, \texttt{PlagiarismDetectionService.java}) to detect cross-document plagiarism across student submissions:
\begin{enumerate}
    \item Submissions $D_1$ and $D_2$ are concatenated with a unique sentinel delimiter:
    \begin{equation}
    T_{\text{combined}} = D_1 \mathbin{\#} D_2 \mathbin{\$}
    \end{equation}
    \item The Suffix Array and Kasai LCP array are evaluated over $T_{\text{combined}}$.
    \item Consecutive entries in $SA$ where one suffix originates in $D_1$ ($pos < |D_1|$) and the other in $D_2$ ($pos > |D_1|$) represent potential shared passages.
    \item Any $LCP[i] \ge \theta$ (threshold $\theta = 40$ characters) is extracted as a plagiarized passage.
\end{enumerate}

\subsubsection{Experimental Demonstration: Alice vs. Bob Plagiarism Case}
In our dataset, student Bob (\texttt{submission\_cs101\_bob.txt}) submitted an essay containing an identical 278-character paragraph copied from Alice (\texttt{submission\_cs101\_alice.txt}):
\begin{quote}
\itshape "For randomized quicksort, partitioning preserves the invariant that all elements preceding the pivot index are strictly less than or equal to the pivot value, while succeeding elements are greater than or equal to the pivot. This structural partition guarantees logarithmic recursion depth under uniform random pivot selection with probability approaching unity as problem dimensions scale."
\end{quote}
Kasai's algorithm detects this excerpt with 100\% recall, calculating an overall document similarity score of 26.33\% and reporting exact line and character boundaries.

\subsection{Screenshot \& Verification Specifications}
\begin{figure}[H]
    \centering
    \fbox{\begin{minipage}{0.85\textwidth}
        \centering
        \vspace{1.2cm}
        \textbf{[PLACEHOLDER: Screenshot 6 - Kasai LCP Plagiarism Report]}\\[0.2cm]
        \textit{Capture the GUI Suffix & Plagiarism Tab comparing Alice vs. Bob submissions.}\\
        \textit{Display: Similarity percentage (26.33\%), detected passage length (278 chars), and side-by-side text preview.}
        \vspace{1.2cm}
    \end{minipage}}
    \caption{Kasai LCP Plagiarism Detection Output in EduTrack Platform}
    \label{fig:screenshot_plagiarism}
\end{figure}

% ------------------------------------------------------------------------------
\section{Suffix Automaton (SAM)}
\subsection{Mathematical Formulation \& Suffix Links}
The Suffix Automaton is the minimal Deterministic Finite Automaton (DFA) recognizing the set of all substrings of text $T$.
\begin{itemize}
    \item \textbf{Endpos Equivalence:} Two substrings $u, v$ belong to the same state if and only if their sets of ending positions in $T$ are identical: $\text{endpos}(u) = \text{endpos}(v)$.
    \item \textbf{Linear Size Bounds:} For a text of length $N \ge 1$:
    \begin{align}
    \text{Number of States} &\le 2N - 1 \nonumber \\
    \text{Number of Transitions} &\le 3N - 4
    \end{align}
    \item \textbf{Suffix Link Tree:} Suffix links connect each state to the state corresponding to its longest suffix with a strictly larger $\text{endpos}$ set.
\end{itemize}

\subsection{Sub-Linear Substring Queries \& Counting}
EduTrack's \texttt{SuffixAutomaton.java} provides:
\begin{enumerate}
    \item \textbf{Substring Presence Test:} Traces transitions character by character in $O(|P|)$ time without binary search.
    \item \textbf{Distinct Substrings Count:} Evaluates the number of distinct substrings in the entire corpus in $O(N)$ dynamic programming over the automaton DAG:
    \begin{equation}
    dp[u] = 1 + \sum_{c \in \Sigma, g(u, c) \ne \text{null}} dp[g(u, c)]
    \end{equation}
    \item \textbf{Longest Common Substring (LCS):} Finds the LCS with another document $D_2$ in $O(|D_2|)$ time by traversing automaton states and backing up along suffix links upon failure.
\end{enumerate}

\subsection{Screenshot \& Verification Specifications}
\begin{figure}[H]
    \centering
    \fbox{\begin{minipage}{0.85\textwidth}
        \centering
        \vspace{1.2cm}
        \textbf{[PLACEHOLDER: Screenshot 7 - Suffix Automaton Substring Query]}\\[0.2cm]
        \textit{Capture the GUI/CLI Suffix Automaton inspection on library document.}\\
        \textit{Display: Total distinct substring count, O(|P|) existence test for "Algorithms", and LCS output.}
        \vspace{1.2cm}
    \end{minipage}}
    \caption{Suffix Automaton State Machine Verification in EduTrack}
    \label{fig:screenshot_sam}
\end{figure}
